% !TEX root = template-shokai-lua.tex
% !TEX recipe = LuaLaTeX (kian-lua)
% ============================================================
%  照会書テンプレート  (kian-v08 LuaLaTeX 版)
%  コンパイル: latexmk -r latexmkrc-lua template-shokai-lua.tex
%              または VS Code の「LuaLaTeX (kian-lua)」レシピ
%  プレースホルダー: 【...】を実際の内容に置き換えてください
%  ※ \SecA{} 等の {} 内に % を書くとコメントになりエラーが出ます
%  見出し番号: \SecZ=番号なし \SecA=第１ \SecB=１ \SecC=⑴ \SecD=ア \SecE=(ア)
%  ※ 見出し（\SecA{} 等）は本文なしで連続させると番号表示が崩れるため、
%    各見出しの直後に【ここに記載】のプレースホルダー本文を入れています
% ============================================================

\documentclass{ltjsarticle}
\usepackage{kian-lua}
% \usepackage{kian-ivs} %IVS 異体字（\ivs{文字}）
% \usepackage{kian-otf} %囲み文字拡張
% \usepackage{luatexja-ruby} %ルビ機能

%  照会書は事件番号欄を付けない（\Idcs は省略）
\Tit{照会書}
\Dat{令和８年　　月　　日}
\Recs{
  [n]{:【照会先機関・会社名】:御中}
}
%\senderhlength{18}{4}{2}{10}% \senderhlength{①}{②}{③}{④}
% ①左端からのインデント（開始位置）② （肩書幅）「弁護士」という肩書を配置するボックスの横幅です。③ （肩書と名前の間隔）「弁護士」と「名前」の間に空けるスペースの幅です。④（名前幅）「【氏名】」という名前を均等割り付け（配置）するボックスの横幅です。
\Sens{
  [0.5]{%
    【事務所名】\\
    弁護士:【氏名】\\
  }
}

% ============================================================

\begin{document}

\maketitle

% ---------------------------------------------------------- %
%  前文
% ---------------------------------------------------------- %

私は，【依頼者氏名】の依頼を受けた代理人弁護士です。
下記のとおり照会しますので，【期日】までにご回答くださいますよう，お願い申し上げます。

% 弁護士法23条の2照会の場合は以下の一文を追加
% なお，本照会は弁護士法第２３条の２に基づくものです。

\vspace{1\zw}

% ---------------------------------------------------------- %
%  照会事項
% ---------------------------------------------------------- %

\SecA{照会事項}

\SecB{}

【ここに記載】

\SecB{}

【ここに記載】

% ---------------------------------------------------------- %
%  照会の目的
% ---------------------------------------------------------- %

\SecA{照会の目的}



% ---------------------------------------------------------- %
%  関係する事件の概要
% ---------------------------------------------------------- %

\SecA{関係する事件の概要}



\vspace{2\zw}

\begin{flushright}
以　上
\end{flushright}

\end{document}
